\documentclass{article}
\begin{document}
\section{Introducción}
Los métodos de inspección tradicionales son costosos y lentos. Se requieren técnicas automatizadas de detección de defectos.
\section{Diseño del Método}
El método propuesto utiliza una arquitectura de dos etapas: la primera realiza preprocesamiento y realce de imágenes, la segunda ejecuta la clasificación de defectos.
La optimización por enjambre de partículas ajusta automáticamente los principales parámetros. La optimización por enjambre de partículas destaca en la exploración global.
La estrategia de aumento de datos incluye recorte aleatorio, jitter de color y deformación elástica.
\section{Análisis de Resultados}
La precisión alcanza 91.3%, superando la línea base en 2.4 puntos. El entrenamiento requiere 50 épocas.
\end{document}
